% SOURCE_AUTHORITY: sga5_fr_workpass.tex, lines 13141--13161 (LF-normalized unit).
% SOURCE_SHA256: 791F4EFFC5E02832D5D77ED03518C8156D6F07E4C8238B03545DB93D883FBB28
Sea para ello $f$ un generador topológico de $\bar\pi$. Se puede suponer que las componentes $M^i$ de $M$ son $\bar\pi$-módulos tales que
\[
H^j(\bar\pi,M^i)=0\quad\text{si }j\ge1.
\]
Existe entonces un triángulo distinguido
\[
\begin{tikzcd}[column sep=huge,row sep=large]
& M^\bullet \arrow[dl] &\\
R\Gamma^{\bar\pi}(M^\bullet) \arrow[rr] && M^\bullet \arrow[ul,"1-f"]
\end{tikzcd}
\tag{7.15}
\]
En efecto, por una parte, es claro que $R\Gamma^{\bar\pi}=\Ker(1-f)$. Por otra parte, $1-f:M^\bullet\to M^\bullet$ es sobreyectivo. En efecto,
\[
H^1(\bar\pi,M)=M^\bullet/(1-f)M^\bullet
\]
(véase para este cálculo [4] p. 197, en un caso sensiblemente análogo), y resulta, por tanto, $M^\bullet=(1-f)M^\bullet$. El triángulo (7.15) da al mismo tiempo $R\Gamma^{\bar\pi}(M^\bullet)\in\Ob(\Delta)$ y
\[
\clK{\Delta}(R\Gamma^{\bar\pi}(M^\bullet))=0,
\]
lo que demuestra 7.14.
\end{prooftext}
